\section{Conclusion}


Quasar represents a fundamental advancement in blockchain consensus design, achieving quantum security without sacrificing performance. Through triple-proof finality combining classical BLS with post-quantum Corona signatures, Q-Chain provides:

\begin{enumerate}
\item \textbf{10\,ms GPU BLS finality} (100\,ms CPU fallback) with triple-proof quantum security
\item \textbf{4{,}761{,}904 TPS} with finality at 1\,B gas block limit
\item \textbf{Quantum resistance} through defense-in-depth
\item \textbf{Modular architecture} supporting various blockchain types
\item \textbf{Physical impossibility} of real-time quantum attacks ($<$10\,ms window)
\end{enumerate}

The narrow $<$10\,ms attack window with GPU BLS makes quantum attacks physically impossible, while the modular consensus stack (Photon, Wave, Nova, Nebula, Prism, Quasar) provides flexibility for different use cases. Q-Chain positions Lux Network at the forefront of blockchain security for the next generation of decentralized applications.

By combining sampling-based consensus, threshold cryptography, and post-quantum signatures, Quasar achieves the seemingly impossible: quantum security with classical-level performance.

\begin{thebibliography}{99}

\bibitem{shor1994}
Shor, P.W. (1994).
\textit{Polynomial-Time Algorithms for Prime Factorization and Discrete Logarithms on a Quantum Computer}.
SIAM Journal on Computing, 26(5), 1484-1509.

\bibitem{nist-pqc}
NIST (2024).
\textit{Post-Quantum Cryptography Standardization}.
National Institute of Standards and Technology.

\bibitem{mosca2018}
Mosca, M. (2018).
\textit{Cybersecurity in an Era with Quantum Computers}.
IEEE Security \& Privacy, 16(5), 38-41.

\bibitem{ntt-corona}
NTT Research (2024).
\textit{Corona: World's First Two-Round Post-Quantum Threshold Signature Scheme}.
Cryptology ePrint Archive.

\bibitem{boneh-bls}
Boneh, D., Lynn, B., \& Shacham, H. (2001).
\textit{Short Signatures from the Weil Pairing}.
Advances in Cryptology---ASIACRYPT 2001, 514-532.

\bibitem{verkle-trees}
Kuszmaul, J. (2019).
\textit{Verkle Trees}.
Ethereum Research.

\bibitem{avalanche-consensus}
Team Rocket (2020).
\textit{Scalable and Probabilistic Leaderless BFT Consensus through Metastability}.
arXiv:1906.08936.

\bibitem{sui-consensus}
Blackshear, S. et al. (2022).
\textit{Narwhal and Tusk: A DAG-based Mempool and Efficient BFT Consensus}.
EuroSys 2022.

\bibitem{lux-benchmarks}
Lux Industries (2026).
\textit{Quasar Consensus --- Measured Benchmarks}.
Apple M1 Max, darwin/arm64, Go 1.26.1.
\url{https://github.com/luxfi/consensus/blob/main/BENCHMARKS.md}.

\end{thebibliography}

\appendix

